\chapter*{Introduction Générale}
\addcontentsline{toc}{chapter}{Introduction Générale}

\noindent La transformation numérique des organisations s'est accompagnée d'une
accélération notable des cycles de développement logiciel. L'adoption généralisée des
méthodologies Agile et des pratiques d'intégration et de déploiement continus
(\textit{CI/CD}) a considérablement réduit le temps séparant l'écriture du code de
sa mise en production. Si cette évolution représente un gain compétitif indéniable,
elle expose simultanément les applications à un risque accru : les vulnérabilités
de sécurité, autrefois détectées lors d'audits ponctuels et planifiés, se propagent
désormais à une vitesse qui dépasse la capacité des équipes à les identifier manuellement.

\medskip
\noindent Dans ce contexte, le paradigme \textbf{DevSecOps} s'est imposé comme une
réponse structurelle à cette problématique. Plutôt que de traiter la sécurité comme
une étape finale du cycle de vie logiciel, le DevSecOps l'intègre dès la conception,
à chaque étape du pipeline de livraison. Cette approche repose sur deux familles
d'analyses complémentaires : l'analyse statique du code source (\textit{SAST — Static
Application Security Testing}), qui identifie les défauts logiques et les
anti-patterns dangereux sans exécuter l'application, et l'analyse dynamique
(\textit{DAST — Dynamic Application Security Testing}), qui teste l'application en
cours d'exécution afin de révéler les failles d'exposition réelles, celles que le
code source seul ne permet pas de détecter.

\medskip
\noindent C'est dans cette perspective qu'a été conçu \textbf{InvisiThreat}, une
plateforme DevSecOps dédiée à la détection automatisée de vulnérabilités sur des
applications web et des services internes. Le projet répond à un besoin concret
observé au sein de l'entreprise Mobelite : l'absence d'un point d'entrée unifié pour
orchestrer les analyses de sécurité, consolider leurs résultats et fournir aux équipes
techniques une vision priorisée et exploitable du risque applicatif.

\medskip
\noindent La solution repose sur une architecture conteneurisée réunissant plusieurs
composants spécialisés. L'analyse statique est assurée par \textbf{Semgrep} et
\textbf{Bandit} pour les sources Python, tandis que les tests dynamiques sont
orchestrés via \textbf{OWASP ZAP} opérant en mode daemon. Les résultats des deux
familles d'analyse sont normalisés, qualifiés selon leur sévérité, puis persistés dans
une base \textbf{PostgreSQL}. Une \textbf{API REST FastAPI} centralise l'ensemble des
opérations et alimente un frontend \textbf{React} offrant tableaux de bord et
fonctionnalités d'export. Une couche de notification temps réel, implémentée via
\textbf{Socket.IO}, informe les utilisateurs dès la finalisation d'un scan ou la
détection d'une vulnérabilité critique.

\medskip
\noindent Au-delà de la détection, InvisiThreat vise la \textbf{traçabilité} : chaque
exécution est rattachée à un projet, une version applicative, un outil de scan et un
horodatage précis. Cette structuration permet de comparer les résultats entre scans
successifs et de mesurer objectivement la réduction du risque au fil des correctifs
appliqués — répondant ainsi à un besoin fondamental des organisations souhaitant
démontrer leur maturité sécurité.

\medskip
\noindent \textbf{Ce rapport s'organise en trois chapitres.} Le premier chapitre
présente le cadre général du projet : contexte organisationnel, problématique de
sécurité, étude des solutions existantes et méthodologie de développement adoptée.
Le deuxième chapitre est consacré à l'étude et la conception du système : spécification
des besoins, identification des acteurs, modélisation UML et description des scénarios
critiques. Le troisième et dernier chapitre détaille l'architecture technique retenue,
les choix d'implémentation des modules clés — pipeline SAST, pipeline DAST,
notifications temps réel, reporting — ainsi que les interfaces réalisées.

\clearpage